% Shared preamble of the Korean reading editions of manuscripts A and B (XeLaTeX).
% Each main.tex loads it with % Shared preamble of the Korean reading editions of manuscripts A and B (XeLaTeX).
% Each main.tex loads it with % Shared preamble of the Korean reading editions of manuscripts A and B (XeLaTeX).
% Each main.tex loads it with % Shared preamble of the Korean reading editions of manuscripts A and B (XeLaTeX).
% Each main.tex loads it with \input{../common/preamble} after \documentclass[11pt]{article}.

\usepackage[a4paper,top=24mm,bottom=26mm,left=25mm,right=25mm,headheight=14pt]{geometry}

% Fonts: Latin in Times/Arial; Hangul in KoPubWorld (with a real bold) where installed,
% then Noto CJK KR, then the Batang and Malgun Gothic that ship with Windows.
\usepackage{fontspec}
\IfFontExistsTF{Times New Roman}{\setmainfont{Times New Roman}}{\setmainfont{TeX Gyre Termes}}
\IfFontExistsTF{Arial}{\setsansfont{Arial}[Scale=0.95]}{\setsansfont{TeX Gyre Heros}[Scale=0.95]}
\IfFontExistsTF{Consolas}{\setmonofont{Consolas}[Scale=0.9]}{}
\usepackage{kotex}
\IfFontExistsTF{KoPubWorldBatang Medium}{%
  \setmainhangulfont{KoPubWorldBatang Medium}[BoldFont={KoPubWorldBatang Bold},AutoFakeSlant=0.15]%
}{\IfFontExistsTF{Noto Serif CJK KR}{%
  \setmainhangulfont{Noto Serif CJK KR}[AutoFakeSlant=0.15]%
}{\setmainhangulfont{Batang}[AutoFakeSlant=0.15]}}
\IfFontExistsTF{KoPubWorldDotum Medium}{%
  \setsanshangulfont{KoPubWorldDotum Medium}[BoldFont={KoPubWorldDotum Bold},AutoFakeSlant=0.15]%
}{\IfFontExistsTF{Noto Sans CJK KR}{%
  \setsanshangulfont{Noto Sans CJK KR}[AutoFakeSlant=0.15]%
}{\setsanshangulfont{Malgun Gothic}[AutoFakeSlant=0.15]}}
\setmonohangulfont{Malgun Gothic}

\linespread{1.3}
\setlength{\parindent}{0pt}
\setlength{\parskip}{0.6em plus 0.2em minus 0.1em}
\setlength{\emergencystretch}{2em}
\raggedbottom

\usepackage{amsmath,amssymb}
\usepackage{array,booktabs,tabularx,multirow}
\usepackage{graphicx}
\usepackage[table]{xcolor}
\usepackage{tikz}
\usetikzlibrary{arrows.meta,positioning,calc}
\usepackage[most]{tcolorbox}
\usepackage{enumitem}
\setlist{itemsep=0.3em,topsep=0.2em,parsep=0.2em,leftmargin=1.7em}
\usepackage[font=small,labelfont={bf,sf},labelsep=period,justification=raggedright,
  singlelinecheck=false,skip=5pt]{caption}

\usepackage{titlesec}
\titleformat{\section}{\sffamily\Large\bfseries}{\thesection}{0.6em}{}
\titleformat{\subsection}{\sffamily\large\bfseries}{\thesubsection}{0.6em}{}
\titleformat{\subsubsection}{\sffamily\normalsize\bfseries}{\thesubsubsection}{0.6em}{}
\titleformat{\paragraph}[runin]{\sffamily\normalsize\bfseries}{}{0pt}{}
\titlespacing*{\section}{0pt}{1.4em}{0.5em}
\titlespacing*{\subsection}{0pt}{1.1em}{0.4em}
\titlespacing*{\subsubsection}{0pt}{0.9em}{0.3em}
\titlespacing*{\paragraph}{0pt}{0.6em}{0.6em}

\usepackage{fancyhdr}
\pagestyle{fancy}
\fancyhf{}
\renewcommand{\headrulewidth}{0.4pt}
\fancyhead[L]{\small\sffamily\ShortTitle}
\fancyhead[R]{\small\sffamily 쉬운 한국어 해설판}
\fancyfoot[C]{\small\thepage}
\newcommand{\ShortTitle}{}

\usepackage{longtable}
\usepackage{cite}
\PassOptionsToPackage{hyphens}{url}
\usepackage[hidelinks]{hyperref}
\urlstyle{same}

\renewcommand{\figurename}{그림}
\renewcommand{\tablename}{표}
\renewcommand{\refname}{참고문헌}
\renewcommand{\contentsname}{차례}

% Boxes. yoji: the gist of a section; puri: an explanation for the reader.
% Neither box is part of the English manuscript; the text outside the boxes is.
\newtcolorbox{yoji}{enhanced,breakable,colback=black!5,colframe=black!5,boxrule=0pt,arc=0pt,
  borderline west={2.5pt}{0pt}{black!50},left=9pt,right=8pt,top=5pt,bottom=5pt,
  before skip=0.5em,after skip=0.9em,fontupper=\small,
  before upper={\textsf{\textbf{요점}}\quad}}
\newtcolorbox{puri}[1][쉽게 풀면]{enhanced,breakable,colback=blue!3,colframe=blue!35!black!45,
  boxrule=0.5pt,arc=1.5pt,left=8pt,right=8pt,top=4pt,bottom=5pt,before skip=0.7em,after skip=0.9em,
  fontupper=\small,title={#1},fonttitle=\sffamily\bfseries\small,coltitle=black!85,
  colbacktitle=blue!9,toptitle=1pt,bottomtitle=1pt}

% Title block: \KoTitle{Korean title}{English title}{authors}{venue and source}
\newcommand{\KoTitle}[4]{%
  \thispagestyle{plain}%
  {\raggedright
  {\sffamily\bfseries\LARGE #1\par}\vspace{0.7em}
  {\large #2\par}\vspace{0.9em}
  {\small #3\par}\vspace{0.3em}
  {\small\color{black!70}#4\par}}
  \vspace{0.6em}\hrule\vspace{0.4em}}

% Numbers written only to illustrate (check_numbers.py skips what this wraps).
\newcommand{\ill}[1]{#1}
\newcommand{\dt}{\Delta\tau}
\newcommand{\Tobs}{T_{\mathrm{obs}}}
% A table cell that wraps.
\newcolumntype{L}[1]{>{\raggedright\arraybackslash}p{#1}}
\newcolumntype{Y}{>{\raggedright\arraybackslash}X}
 after \documentclass[11pt]{article}.

\usepackage[a4paper,top=24mm,bottom=26mm,left=25mm,right=25mm,headheight=14pt]{geometry}

% Fonts: Latin in Times/Arial; Hangul in KoPubWorld (with a real bold) where installed,
% then Noto CJK KR, then the Batang and Malgun Gothic that ship with Windows.
\usepackage{fontspec}
\IfFontExistsTF{Times New Roman}{\setmainfont{Times New Roman}}{\setmainfont{TeX Gyre Termes}}
\IfFontExistsTF{Arial}{\setsansfont{Arial}[Scale=0.95]}{\setsansfont{TeX Gyre Heros}[Scale=0.95]}
\IfFontExistsTF{Consolas}{\setmonofont{Consolas}[Scale=0.9]}{}
\usepackage{kotex}
\IfFontExistsTF{KoPubWorldBatang Medium}{%
  \setmainhangulfont{KoPubWorldBatang Medium}[BoldFont={KoPubWorldBatang Bold},AutoFakeSlant=0.15]%
}{\IfFontExistsTF{Noto Serif CJK KR}{%
  \setmainhangulfont{Noto Serif CJK KR}[AutoFakeSlant=0.15]%
}{\setmainhangulfont{Batang}[AutoFakeSlant=0.15]}}
\IfFontExistsTF{KoPubWorldDotum Medium}{%
  \setsanshangulfont{KoPubWorldDotum Medium}[BoldFont={KoPubWorldDotum Bold},AutoFakeSlant=0.15]%
}{\IfFontExistsTF{Noto Sans CJK KR}{%
  \setsanshangulfont{Noto Sans CJK KR}[AutoFakeSlant=0.15]%
}{\setsanshangulfont{Malgun Gothic}[AutoFakeSlant=0.15]}}
\setmonohangulfont{Malgun Gothic}

\linespread{1.3}
\setlength{\parindent}{0pt}
\setlength{\parskip}{0.6em plus 0.2em minus 0.1em}
\setlength{\emergencystretch}{2em}
\raggedbottom

\usepackage{amsmath,amssymb}
\usepackage{array,booktabs,tabularx,multirow}
\usepackage{graphicx}
\usepackage[table]{xcolor}
\usepackage{tikz}
\usetikzlibrary{arrows.meta,positioning,calc}
\usepackage[most]{tcolorbox}
\usepackage{enumitem}
\setlist{itemsep=0.3em,topsep=0.2em,parsep=0.2em,leftmargin=1.7em}
\usepackage[font=small,labelfont={bf,sf},labelsep=period,justification=raggedright,
  singlelinecheck=false,skip=5pt]{caption}

\usepackage{titlesec}
\titleformat{\section}{\sffamily\Large\bfseries}{\thesection}{0.6em}{}
\titleformat{\subsection}{\sffamily\large\bfseries}{\thesubsection}{0.6em}{}
\titleformat{\subsubsection}{\sffamily\normalsize\bfseries}{\thesubsubsection}{0.6em}{}
\titleformat{\paragraph}[runin]{\sffamily\normalsize\bfseries}{}{0pt}{}
\titlespacing*{\section}{0pt}{1.4em}{0.5em}
\titlespacing*{\subsection}{0pt}{1.1em}{0.4em}
\titlespacing*{\subsubsection}{0pt}{0.9em}{0.3em}
\titlespacing*{\paragraph}{0pt}{0.6em}{0.6em}

\usepackage{fancyhdr}
\pagestyle{fancy}
\fancyhf{}
\renewcommand{\headrulewidth}{0.4pt}
\fancyhead[L]{\small\sffamily\ShortTitle}
\fancyhead[R]{\small\sffamily 쉬운 한국어 해설판}
\fancyfoot[C]{\small\thepage}
\newcommand{\ShortTitle}{}

\usepackage{longtable}
\usepackage{cite}
\PassOptionsToPackage{hyphens}{url}
\usepackage[hidelinks]{hyperref}
\urlstyle{same}

\renewcommand{\figurename}{그림}
\renewcommand{\tablename}{표}
\renewcommand{\refname}{참고문헌}
\renewcommand{\contentsname}{차례}

% Boxes. yoji: the gist of a section; puri: an explanation for the reader.
% Neither box is part of the English manuscript; the text outside the boxes is.
\newtcolorbox{yoji}{enhanced,breakable,colback=black!5,colframe=black!5,boxrule=0pt,arc=0pt,
  borderline west={2.5pt}{0pt}{black!50},left=9pt,right=8pt,top=5pt,bottom=5pt,
  before skip=0.5em,after skip=0.9em,fontupper=\small,
  before upper={\textsf{\textbf{요점}}\quad}}
\newtcolorbox{puri}[1][쉽게 풀면]{enhanced,breakable,colback=blue!3,colframe=blue!35!black!45,
  boxrule=0.5pt,arc=1.5pt,left=8pt,right=8pt,top=4pt,bottom=5pt,before skip=0.7em,after skip=0.9em,
  fontupper=\small,title={#1},fonttitle=\sffamily\bfseries\small,coltitle=black!85,
  colbacktitle=blue!9,toptitle=1pt,bottomtitle=1pt}

% Title block: \KoTitle{Korean title}{English title}{authors}{venue and source}
\newcommand{\KoTitle}[4]{%
  \thispagestyle{plain}%
  {\raggedright
  {\sffamily\bfseries\LARGE #1\par}\vspace{0.7em}
  {\large #2\par}\vspace{0.9em}
  {\small #3\par}\vspace{0.3em}
  {\small\color{black!70}#4\par}}
  \vspace{0.6em}\hrule\vspace{0.4em}}

% Numbers written only to illustrate (check_numbers.py skips what this wraps).
\newcommand{\ill}[1]{#1}
\newcommand{\dt}{\Delta\tau}
\newcommand{\Tobs}{T_{\mathrm{obs}}}
% A table cell that wraps.
\newcolumntype{L}[1]{>{\raggedright\arraybackslash}p{#1}}
\newcolumntype{Y}{>{\raggedright\arraybackslash}X}
 after \documentclass[11pt]{article}.

\usepackage[a4paper,top=24mm,bottom=26mm,left=25mm,right=25mm,headheight=14pt]{geometry}

% Fonts: Latin in Times/Arial; Hangul in KoPubWorld (with a real bold) where installed,
% then Noto CJK KR, then the Batang and Malgun Gothic that ship with Windows.
\usepackage{fontspec}
\IfFontExistsTF{Times New Roman}{\setmainfont{Times New Roman}}{\setmainfont{TeX Gyre Termes}}
\IfFontExistsTF{Arial}{\setsansfont{Arial}[Scale=0.95]}{\setsansfont{TeX Gyre Heros}[Scale=0.95]}
\IfFontExistsTF{Consolas}{\setmonofont{Consolas}[Scale=0.9]}{}
\usepackage{kotex}
\IfFontExistsTF{KoPubWorldBatang Medium}{%
  \setmainhangulfont{KoPubWorldBatang Medium}[BoldFont={KoPubWorldBatang Bold},AutoFakeSlant=0.15]%
}{\IfFontExistsTF{Noto Serif CJK KR}{%
  \setmainhangulfont{Noto Serif CJK KR}[AutoFakeSlant=0.15]%
}{\setmainhangulfont{Batang}[AutoFakeSlant=0.15]}}
\IfFontExistsTF{KoPubWorldDotum Medium}{%
  \setsanshangulfont{KoPubWorldDotum Medium}[BoldFont={KoPubWorldDotum Bold},AutoFakeSlant=0.15]%
}{\IfFontExistsTF{Noto Sans CJK KR}{%
  \setsanshangulfont{Noto Sans CJK KR}[AutoFakeSlant=0.15]%
}{\setsanshangulfont{Malgun Gothic}[AutoFakeSlant=0.15]}}
\setmonohangulfont{Malgun Gothic}

\linespread{1.3}
\setlength{\parindent}{0pt}
\setlength{\parskip}{0.6em plus 0.2em minus 0.1em}
\setlength{\emergencystretch}{2em}
\raggedbottom

\usepackage{amsmath,amssymb}
\usepackage{array,booktabs,tabularx,multirow}
\usepackage{graphicx}
\usepackage[table]{xcolor}
\usepackage{tikz}
\usetikzlibrary{arrows.meta,positioning,calc}
\usepackage[most]{tcolorbox}
\usepackage{enumitem}
\setlist{itemsep=0.3em,topsep=0.2em,parsep=0.2em,leftmargin=1.7em}
\usepackage[font=small,labelfont={bf,sf},labelsep=period,justification=raggedright,
  singlelinecheck=false,skip=5pt]{caption}

\usepackage{titlesec}
\titleformat{\section}{\sffamily\Large\bfseries}{\thesection}{0.6em}{}
\titleformat{\subsection}{\sffamily\large\bfseries}{\thesubsection}{0.6em}{}
\titleformat{\subsubsection}{\sffamily\normalsize\bfseries}{\thesubsubsection}{0.6em}{}
\titleformat{\paragraph}[runin]{\sffamily\normalsize\bfseries}{}{0pt}{}
\titlespacing*{\section}{0pt}{1.4em}{0.5em}
\titlespacing*{\subsection}{0pt}{1.1em}{0.4em}
\titlespacing*{\subsubsection}{0pt}{0.9em}{0.3em}
\titlespacing*{\paragraph}{0pt}{0.6em}{0.6em}

\usepackage{fancyhdr}
\pagestyle{fancy}
\fancyhf{}
\renewcommand{\headrulewidth}{0.4pt}
\fancyhead[L]{\small\sffamily\ShortTitle}
\fancyhead[R]{\small\sffamily 쉬운 한국어 해설판}
\fancyfoot[C]{\small\thepage}
\newcommand{\ShortTitle}{}

\usepackage{longtable}
\usepackage{cite}
\PassOptionsToPackage{hyphens}{url}
\usepackage[hidelinks]{hyperref}
\urlstyle{same}

\renewcommand{\figurename}{그림}
\renewcommand{\tablename}{표}
\renewcommand{\refname}{참고문헌}
\renewcommand{\contentsname}{차례}

% Boxes. yoji: the gist of a section; puri: an explanation for the reader.
% Neither box is part of the English manuscript; the text outside the boxes is.
\newtcolorbox{yoji}{enhanced,breakable,colback=black!5,colframe=black!5,boxrule=0pt,arc=0pt,
  borderline west={2.5pt}{0pt}{black!50},left=9pt,right=8pt,top=5pt,bottom=5pt,
  before skip=0.5em,after skip=0.9em,fontupper=\small,
  before upper={\textsf{\textbf{요점}}\quad}}
\newtcolorbox{puri}[1][쉽게 풀면]{enhanced,breakable,colback=blue!3,colframe=blue!35!black!45,
  boxrule=0.5pt,arc=1.5pt,left=8pt,right=8pt,top=4pt,bottom=5pt,before skip=0.7em,after skip=0.9em,
  fontupper=\small,title={#1},fonttitle=\sffamily\bfseries\small,coltitle=black!85,
  colbacktitle=blue!9,toptitle=1pt,bottomtitle=1pt}

% Title block: \KoTitle{Korean title}{English title}{authors}{venue and source}
\newcommand{\KoTitle}[4]{%
  \thispagestyle{plain}%
  {\raggedright
  {\sffamily\bfseries\LARGE #1\par}\vspace{0.7em}
  {\large #2\par}\vspace{0.9em}
  {\small #3\par}\vspace{0.3em}
  {\small\color{black!70}#4\par}}
  \vspace{0.6em}\hrule\vspace{0.4em}}

% Numbers written only to illustrate (check_numbers.py skips what this wraps).
\newcommand{\ill}[1]{#1}
\newcommand{\dt}{\Delta\tau}
\newcommand{\Tobs}{T_{\mathrm{obs}}}
% A table cell that wraps.
\newcolumntype{L}[1]{>{\raggedright\arraybackslash}p{#1}}
\newcolumntype{Y}{>{\raggedright\arraybackslash}X}
 after \documentclass[11pt]{article}.

\usepackage[a4paper,top=24mm,bottom=26mm,left=25mm,right=25mm,headheight=14pt]{geometry}

% Fonts: Latin in Times/Arial; Hangul in KoPubWorld (with a real bold) where installed,
% then Noto CJK KR, then the Batang and Malgun Gothic that ship with Windows.
\usepackage{fontspec}
\IfFontExistsTF{Times New Roman}{\setmainfont{Times New Roman}}{\setmainfont{TeX Gyre Termes}}
\IfFontExistsTF{Arial}{\setsansfont{Arial}[Scale=0.95]}{\setsansfont{TeX Gyre Heros}[Scale=0.95]}
\IfFontExistsTF{Consolas}{\setmonofont{Consolas}[Scale=0.9]}{}
\usepackage{kotex}
\IfFontExistsTF{KoPubWorldBatang Medium}{%
  \setmainhangulfont{KoPubWorldBatang Medium}[BoldFont={KoPubWorldBatang Bold},AutoFakeSlant=0.15]%
}{\IfFontExistsTF{Noto Serif CJK KR}{%
  \setmainhangulfont{Noto Serif CJK KR}[AutoFakeSlant=0.15]%
}{\setmainhangulfont{Batang}[AutoFakeSlant=0.15]}}
\IfFontExistsTF{KoPubWorldDotum Medium}{%
  \setsanshangulfont{KoPubWorldDotum Medium}[BoldFont={KoPubWorldDotum Bold},AutoFakeSlant=0.15]%
}{\IfFontExistsTF{Noto Sans CJK KR}{%
  \setsanshangulfont{Noto Sans CJK KR}[AutoFakeSlant=0.15]%
}{\setsanshangulfont{Malgun Gothic}[AutoFakeSlant=0.15]}}
\setmonohangulfont{Malgun Gothic}

\linespread{1.3}
\setlength{\parindent}{0pt}
\setlength{\parskip}{0.6em plus 0.2em minus 0.1em}
\setlength{\emergencystretch}{2em}
\raggedbottom

\usepackage{amsmath,amssymb}
\usepackage{array,booktabs,tabularx,multirow}
\usepackage{graphicx}
\usepackage[table]{xcolor}
\usepackage{tikz}
\usetikzlibrary{arrows.meta,positioning,calc}
\usepackage[most]{tcolorbox}
\usepackage{enumitem}
\setlist{itemsep=0.3em,topsep=0.2em,parsep=0.2em,leftmargin=1.7em}
\usepackage[font=small,labelfont={bf,sf},labelsep=period,justification=raggedright,
  singlelinecheck=false,skip=5pt]{caption}

\usepackage{titlesec}
\titleformat{\section}{\sffamily\Large\bfseries}{\thesection}{0.6em}{}
\titleformat{\subsection}{\sffamily\large\bfseries}{\thesubsection}{0.6em}{}
\titleformat{\subsubsection}{\sffamily\normalsize\bfseries}{\thesubsubsection}{0.6em}{}
\titleformat{\paragraph}[runin]{\sffamily\normalsize\bfseries}{}{0pt}{}
\titlespacing*{\section}{0pt}{1.4em}{0.5em}
\titlespacing*{\subsection}{0pt}{1.1em}{0.4em}
\titlespacing*{\subsubsection}{0pt}{0.9em}{0.3em}
\titlespacing*{\paragraph}{0pt}{0.6em}{0.6em}

\usepackage{fancyhdr}
\pagestyle{fancy}
\fancyhf{}
\renewcommand{\headrulewidth}{0.4pt}
\fancyhead[L]{\small\sffamily\ShortTitle}
\fancyhead[R]{\small\sffamily 쉬운 한국어 해설판}
\fancyfoot[C]{\small\thepage}
\newcommand{\ShortTitle}{}

\usepackage{longtable}
\usepackage{cite}
\PassOptionsToPackage{hyphens}{url}
\usepackage[hidelinks]{hyperref}
\urlstyle{same}

\renewcommand{\figurename}{그림}
\renewcommand{\tablename}{표}
\renewcommand{\refname}{참고문헌}
\renewcommand{\contentsname}{차례}

% Boxes. yoji: the gist of a section; puri: an explanation for the reader.
% Neither box is part of the English manuscript; the text outside the boxes is.
\newtcolorbox{yoji}{enhanced,breakable,colback=black!5,colframe=black!5,boxrule=0pt,arc=0pt,
  borderline west={2.5pt}{0pt}{black!50},left=9pt,right=8pt,top=5pt,bottom=5pt,
  before skip=0.5em,after skip=0.9em,fontupper=\small,
  before upper={\textsf{\textbf{요점}}\quad}}
\newtcolorbox{puri}[1][쉽게 풀면]{enhanced,breakable,colback=blue!3,colframe=blue!35!black!45,
  boxrule=0.5pt,arc=1.5pt,left=8pt,right=8pt,top=4pt,bottom=5pt,before skip=0.7em,after skip=0.9em,
  fontupper=\small,title={#1},fonttitle=\sffamily\bfseries\small,coltitle=black!85,
  colbacktitle=blue!9,toptitle=1pt,bottomtitle=1pt}

% Title block: \KoTitle{Korean title}{English title}{authors}{venue and source}
\newcommand{\KoTitle}[4]{%
  \thispagestyle{plain}%
  {\raggedright
  {\sffamily\bfseries\LARGE #1\par}\vspace{0.7em}
  {\large #2\par}\vspace{0.9em}
  {\small #3\par}\vspace{0.3em}
  {\small\color{black!70}#4\par}}
  \vspace{0.6em}\hrule\vspace{0.4em}}

% Numbers written only to illustrate (check_numbers.py skips what this wraps).
\newcommand{\ill}[1]{#1}
\newcommand{\dt}{\Delta\tau}
\newcommand{\Tobs}{T_{\mathrm{obs}}}
% A table cell that wraps.
\newcolumntype{L}[1]{>{\raggedright\arraybackslash}p{#1}}
\newcolumntype{Y}{>{\raggedright\arraybackslash}X}
